\scanpage{94}
\mspage{46}
On peut généraliser encore un peu ces constructions, en \deleted{choisissant une \unsure{partie}} \inserted{remplaçant $\unsure{X'|\widehat X}$ par} \inserted{une réunion de composantes connexes} \deleted{de composantes irréductibles}, contenant $A'$ et $I'$. Cela revient à prendre un sous-graphe \uline{quelconque} $G'$ de $G$ (n'ayant pas $\text{néc}^{\text{t}}$ \deleted{\unsure{les}} $=$ ens. de sommets, mais un ens. de sommets $S'\subset S$), [défini donc par la donnée d'une partie $S'$ de $S$, d'une partie $\widetilde A'$ de $\widetilde A(G)$ \unsure{st.ble} par $\sigma_{\widetilde A}$ et telle que $o_G:\widetilde A\longrightarrow S$ applique $\widetilde A'$ dans $S'$], plus une partie $I'$ de $I\amalg\deleted{\widetilde A}\widetilde J$ (où $\widetilde J=\widetilde A\setminus\widetilde A'$), appliquée dans $S'$ par l'\deleted{application origine}\inserted{application}
\[
I\amalg\widetilde J\longrightarrow S\quad\text{induite par }I\longrightarrow S\text{ et par }\widetilde A\xrightarrow{o_G}S.
\]

Un cas particulièrement important est celui où $G'$ est un sous-graphe ponctuel, défini par un $\alpha\in S$, \inserted{donc $\widetilde J=\widetilde A$,} et où on prend
\[
I'=I_\alpha\amalg\widetilde A_\alpha=\text{image inverse de }\{\alpha\}\text{ par }I\amalg\widetilde A\longrightarrow S.
\]

Si on veut exprimer ces constructions en termes d'une notion ad-hoc de
\margin{93}
